\chapter{Periodic Systems and Bloch's Theorem in One Dimension}
\index{Bloch theorem!one-dimensional|textbf}
\index{periodic boundary condition!one-dimensional|textbf}
\index{reciprocal lattice!one-dimensional|textbf}
\index{Brillouin zone!one-dimensional|textbf}

\label{ch:training-bloch-theorem-one-dimensional}

\chapterstatus{pedagogical bridge from particle-in-a-box models to periodic
quantum systems. The derivations and examples are illustrative mathematics;
this chapter reports no material calculation, numerical convergence claim,
scientific validation, or uncertainty quantification.}

\section{Learning goals}

After working through this chapter, a reader should be able to

\begin{enumerate}
  \item distinguish a Dirichlet box from a periodic system and a fixed Bloch
        sector;
  \item derive the one-dimensional Bloch condition from translation symmetry;
  \item distinguish the Bloch wave number $k$, the primitive reciprocal
        magnitude $b$, and a reciprocal-lattice vector $G_m=mb$;
  \item interpret $k$ modulo reciprocal-lattice translations and choose a
        first-Brillouin-zone representative;
  \item factor a Bloch wave into a phase and a cell-periodic function;
  \item construct the cell-periodic fiber operator for a periodic Schr\"odinger
        Hamiltonian; and
  \item identify the plane-wave momenta $k+G_m$ that occur inside one Bloch
        fiber.
\end{enumerate}

Chapters~\ref{ch:training-quantum-operators} and
\ref{ch:training-particle-box-higher-dimensions} used finite intervals and
rectangular domains with Dirichlet boundary conditions. We now change the
physical model. A periodic crystal is not obtained by relabeling the walls of a
box. It is obtained by declaring translation symmetry and then resolving the
infinite system into sectors characterized by the phase acquired under one
primitive translation.

This one-dimensional construction introduces Bloch's theorem before the
notation is generalized to two and three dimensions. It also avoids a collision
with the Gram metric used in the next chapter: $b=2\pi/a$ is the primitive
reciprocal magnitude in one dimension, and $G_m=mb$ is a reciprocal-lattice
vector indexed by the integer $m$.

\section{Three boundary models that must not be conflated}
\index{Dirichlet boundary condition}
\index{Born--von Karman boundary condition}
\index{Bloch boundary condition}

Let $a>0$ be a length. Three superficially similar interval calculations
represent different operators.

\begin{description}
  \item[Dirichlet box.] A wavefunction on $[0,a]$ obeys
    \begin{equation}
      \psi(0)=\psi(a)=0.
    \end{equation}
    The endpoints are walls, and translation by $a$ is not a symmetry of the
    operator.
  \item[Periodic ring or Born--von Karman cell.] A wavefunction obeys
    \begin{equation}
      \psi(x+a)=\psi(x).
    \end{equation}
    The endpoints are two coordinate representatives of one periodic point.
  \item[Bloch sector.] A wavefunction obeys
    \begin{equation}
      \psi_k(x+a)=e^{ika}\psi_k(x).
      \label{eq:training-bloch-1d-boundary-condition}
    \end{equation}
    The endpoints differ by a phase fixed by the sector label $k$.
\end{description}

The periodic ring is the special Bloch sector for which $e^{ika}=1$. A generic
Bloch sector is quasiperiodic rather than periodic. None of these conditions is
a numerical approximation to either of the others; each is part of a different
operator domain.

\section{Translation symmetry and Bloch's theorem}
\index{translation operator}
\index{quasimomentum}

Let the infinite one-dimensional Hamiltonian be
\begin{equation}
  \widehat H
  =-\frac{\hbar^2}{2m_0}\frac{\mathrm d^2}{\mathrm d x^2}+V(x),
  \qquad
  V(x+a)=V(x),
  \label{eq:training-bloch-1d-parent}
\end{equation}
where $m_0$ is the particle mass. Define translation by one primitive period as
\begin{equation}
  (\widehat T_a\psi)(x)=\psi(x+a).
\end{equation}
The periodicity of $V$ implies
\begin{equation}
  [\widehat H,\widehat T_a]=0.
\end{equation}
Translation is unitary, so any translation eigenvalue has unit magnitude. Write
such a phase as $e^{ika}$. On the infinite line, a nonzero function whose
magnitude repeats in every cell is not square integrable; Bloch waves are
therefore generalized states represented and normalized on one cell. They
become ordinary normalizable eigenstates after a finite Born--von Karman
supercell is imposed. In either formulation, simultaneous resolution by energy
and translation gives
\begin{equation}
  \widehat T_a\psi_{nk}=e^{ika}\psi_{nk},
\end{equation}
which is precisely Equation~\eqref{eq:training-bloch-1d-boundary-condition}.
The label $n$ distinguishes energy bands within the same $k$ sector.

\begin{definition}[One-dimensional Bloch sector]
The Bloch sector labelled by $k$ consists of functions square integrable on one
cell whose quasiperiodic extensions satisfy
\begin{equation}
  \psi_k(x+\ell a)=e^{ik\ell a}\psi_k(x),
  \qquad \ell\in\mathbb Z.
  \label{eq:training-bloch-1d-sector}
\end{equation}
The corresponding Bloch-fiber operator is the restriction of the periodic
Hamiltonian to that sector.
\end{definition}

This is the operational content of Bloch's theorem needed here: generalized
states of an infinite translation-invariant periodic Hamiltonian, or ordinary
eigenstates of its finite Born--von Karman counterpart, can be chosen to have a
definite translation phase~\cite{kittel2004,simon2013}. It does not say that the
wavefunction itself is periodic.

\section{The reciprocal lattice and the first Brillouin zone}
\index{reciprocal-lattice vector!one-dimensional}
\index{reduced momentum}

The primitive reciprocal magnitude is
\begin{equation}
  b=\frac{2\pi}{a}.
  \label{eq:training-geometry-one-dimensional-reciprocal}
\end{equation}
Every reciprocal-lattice vector is
\begin{equation}
  G_m=mb,
  \qquad m\in\mathbb Z.
  \label{eq:training-bloch-1d-reciprocal-vector}
\end{equation}
It obeys
\begin{equation}
  e^{iG_m\ell a}=e^{i2\pi m\ell}=1.
\end{equation}
Consequently, $k$ and $k+G_m$ define the same translation phase:
\begin{equation}
  e^{i(k+G_m)\ell a}=e^{ik\ell a}.
\end{equation}
They label the same reciprocal class,
\begin{equation}
  k\sim k+G_m.
  \label{eq:training-bloch-1d-equivalence}
\end{equation}

Define the reduced, dimensionless coordinate
\begin{equation}
  \kappa=\frac{k}{b}.
  \label{eq:training-geometry-reduced-momentum}
\end{equation}
A convenient first-Brillouin-zone representative is
\begin{equation}
  -\frac12\leq\kappa<\frac12,
  \qquad\text{equivalently}\qquad
  -\frac b2\leq k<\frac b2.
  \label{eq:training-geometry-first-brillouin-zone}
\end{equation}
The half-open interval avoids storing both endpoints of the same reciprocal
class. A closed plotting path may display both endpoints, but a complete
periodic mesh should not count them as distinct samples.

\section{Bloch factorization and the cell-periodic space}
\index{cell-periodic function}

Define
\begin{equation}
  u_{nk}(x)=e^{-ikx}\psi_{nk}(x).
\end{equation}
Using the Bloch condition gives
\begin{align}
  u_{nk}(x+a)
  &=e^{-ik(x+a)}\psi_{nk}(x+a)\\
  &=e^{-ikx}e^{-ika}e^{ika}\psi_{nk}(x)\\
  &=u_{nk}(x).
\end{align}
Thus
\begin{equation}
  \psi_{nk}(x)=e^{ikx}u_{nk}(x),
  \qquad
  u_{nk}(x+a)=u_{nk}(x).
  \label{eq:training-geometry-bloch-factorization}
\end{equation}
The phase factor carries the translation character; the function $u_{nk}$ is
periodic on one primitive cell.

The factorization is not unique under reciprocal relabeling. Replacing $k$ by
$k+G_m$ and $u_{nk}$ by $e^{-iG_mx}u_{nk}$ leaves the physical wavefunction
unchanged. This is why a stored $k$ value must be interpreted together with its
chosen Brillouin-zone representative and basis convention.

\section{Cell-periodic fiber operators}
\index{Bloch-fiber operator!one-dimensional}

Work in the fixed periodic space
\begin{equation}
  \mathcal H_{\mathrm{per}}=L^2_{\mathrm{per}}([0,a]).
\end{equation}
Conjugating the parent Hamiltonian by the Bloch phase defines
\begin{equation}
  \widehat H(k)
  =e^{-ikx}\widehat H e^{ikx}
  \quad\text{on }\mathcal H_{\mathrm{per}}.
  \label{eq:training-geometry-fiber-definition}
\end{equation}
For Equation~\eqref{eq:training-bloch-1d-parent},
\begin{equation}
  \widehat H(k)
  =\frac{\hbar^2}{2m_0}
   \left(-i\frac{\mathrm d}{\mathrm d x}+k\right)^2+V(x).
  \label{eq:training-geometry-schrodinger-fiber}
\end{equation}
The $k$ dependence belongs to the represented derivative and boundary sector;
it is not an additional potential.

At the operator level, the Floquet--Bloch transformation assembles the
continuum of fibers as a direct integral,
\begin{equation}
  \widehat H
  \simeq
  \int_{[-b/2,b/2)}^{\oplus}
    \widehat H(k)\,\frac{a\,\mathrm d k}{2\pi}.
  \label{eq:training-geometry-direct-integral}
\end{equation}
Here $\simeq$ means unitary equivalence and $\oplus$ indicates one fiber space
for each reciprocal class~\cite{kuchment2016}. This formula is an optional
operator-level summary; the later finite calculations require only the
individual fibers and their declared $k$ samples.

The primitive reciprocal magnitude defines the kinetic-energy scale
\begin{equation}
  E_G=\frac{\hbar^2b^2}{2m_0}.
  \label{eq:training-geometry-reciprocal-energy}
\end{equation}
The conventional label $E_G$ is retained for compatibility with later
calculation records, but its defining reciprocal magnitude is $b$. It does not
introduce a standalone reciprocal magnitude called $G$.

\section{Plane waves inside one Bloch fiber}
\index{plane-wave basis!one-dimensional Bloch fiber}

A periodic function on the cell has Fourier basis
\begin{equation}
  \phi_m(x)=\frac{1}{\sqrt a}e^{iG_mx}
  =\frac{1}{\sqrt a}e^{imbx},
  \qquad m\in\mathbb Z.
  \label{eq:training-bloch-1d-cell-plane-wave}
\end{equation}
Multiplying by the Bloch phase gives the corresponding full-space factor
\begin{equation}
  e^{ikx}\phi_m(x)
  =\frac{1}{\sqrt a}e^{i(k+G_m)x}.
\end{equation}
Thus $k$ labels the Bloch sector while $G_m$ labels a basis function inside
that sector. Confusing these roles obscures band folding and reciprocal-basis
ordering.

For the free particle, the fiber energies are
\begin{equation}
  E_m^{(0)}(k)
  =\frac{\hbar^2(k+G_m)^2}{2m_0}.
  \label{eq:training-bloch-1d-free-folded-bands}
\end{equation}
As $k$ is restricted to the first Brillouin zone, the different integers $m$
produce the familiar folded parabolas. A periodic potential couples these
reciprocal basis functions according to its Fourier coefficients. The cosine
parent used later has only the harmonics $G_{+1}=b$ and $G_{-1}=-b$, so it
couples neighboring reciprocal indices.

\section{Finite periodic samples and Born--von Karman quantization}

A finite ring of $N$ primitive cells has length $L=Na$. Imposing
\begin{equation}
  \psi(x+L)=\psi(x)
\end{equation}
requires
\begin{equation}
  e^{ikL}=1,
  \qquad
  k_j=\frac{2\pi j}{Na},
  \qquad j\in\mathbb Z.
\end{equation}
Modulo $b=2\pi/a$, this produces $N$ distinct reciprocal classes. A finite
$k$ mesh is therefore a declared sampling of the Brillouin zone. Refining that
mesh changes the sampling resolution; it does not change Bloch's theorem or
turn a finite sample into the continuum of all $k$ values.

\documentcallout{Notebook to do}{Create
\path{examples/tutorials/research-monograph/foundations/bloch_theorem_one_dimensional.ipynb}
with explicit imports, explanatory inline comments, Dirichlet/periodic/Bloch
boundary comparisons, reciprocal equivalence, cell-periodic factorization,
folded free-particle bands, and a finite Born--von Karman mesh. Keep all examples
synthetic and the notebook unexecuted in the repository.}

\section{A comparison table}

\begin{center}
\begin{tabular}{@{}lll@{}}
\toprule
Model & Boundary statement & Natural modes \\
\midrule
Dirichlet box & $\psi(0)=\psi(a)=0$ & $\sin(n\pi x/a)$ \\
Periodic cell & $\psi(x+a)=\psi(x)$ & $e^{iG_mx}$ \\
Bloch sector & $\psi_k(x+a)=e^{ika}\psi_k(x)$ & $e^{i(k+G_m)x}$ \\
\bottomrule
\end{tabular}
\end{center}

The table compares operator domains and mode labels. It does not identify the
hard-wall box with a periodic crystal. The box provides a familiar eigenproblem;
translation symmetry supplies the new physical structure.

\section{Common mistakes}

\begin{enumerate}
  \item Treating the Dirichlet interval as a primitive periodic cell without
        changing the boundary conditions and operator domain.
  \item Calling $\psi_{nk}$ periodic when only $u_{nk}$ is cell periodic.
  \item Confusing the Bloch wave number $k$ with the reciprocal-lattice vector
        $G_m$.
  \item Counting both endpoints of a half-open Brillouin-zone mesh as distinct
        reciprocal classes.
  \item Treating a finite Born--von Karman mesh as the continuum of all fibers.
  \item Calling a plane-wave cutoff a new physical Bloch sector rather than a
        finite basis selection inside one sector.
\end{enumerate}

\section{Exercises}

\begin{enumerate}
  \item In a finite Born--von Karman supercell, start from
        $[\widehat H,\widehat T_a]=0$ and explain why a nondegenerate energy
        eigenstate may be chosen as a translation eigenstate.
  \item Verify directly that $k$ and $k+G_m$ give the same phase under every
        translation by $\ell a$.
  \item Derive the periodicity of $u_{nk}$ from the Bloch condition.
  \item Derive Equation~\eqref{eq:training-geometry-schrodinger-fiber} by
        applying the parent Hamiltonian to $e^{ikx}u(x)$.
  \item For a ring of four cells, list four inequivalent $k$ representatives in
        the half-open first Brillouin zone.
  \item Sketch the three folded free-particle branches with $m=-1,0,1$ over
        $-b/2\leq k<b/2$ and identify their crossings.
  \item Explain why changing a Dirichlet wall to a Bloch phase is a change of
        operator domain rather than a gauge transformation inside one fixed
        problem.
\end{enumerate}

\section{Evidence boundary and next use}

Bloch's theorem follows here from the declared translation symmetry of the
one-dimensional parent. A finite basis or finite $k$ mesh is a separate
representation choice. Neither the theorem nor the folded free-particle
example establishes convergence of a numerical discretization or adequacy for
a material.

Chapter~\ref{ch:training-periodic-geometry} next generalizes direct and
reciprocal lattices, Brillouin-zone coordinates, and the Bloch condition to two
and three dimensions, including a two-dimensional plane embedded in ambient
three-space.
